% =====================================================================
%  Literature Review chapter -- AI-Assisted Rooftop Solar Potential
%  Estimation and Financial Analysis System.
%
%  Like methodology.tex, this file is a FRAGMENT, intended to be
%  \input{} from a master thesis file (e.g. bachelor.tex) that
%  provides the document class, preamble, \begin{document},
%  \bibliography{references}, etc.
%
%  Required packages (already listed in methodology.tex preamble):
%      booktabs, tabularx, hyperref, natbib (round, authoryear)
% =====================================================================

\providecommand{\degC}{\,^{\circ}\mathrm{C}}
\providecommand{\COtwo}{\mathrm{CO_{2}}}

% =====================================================================
\chapter{Literature Review}
\label{ch:litreview}
% Replace \chapter with \section if your thesis class is `article`.
% =====================================================================

\noindent\textbf{Status:} Chapter~2 deliverable --- academic, LaTeX-ready.\\
\textbf{Pairs with:} \texttt{methodology.tex} (Chapter~3),
\texttt{references.bib}. All citations below resolve against
\texttt{references.bib}.

\bigskip

% =====================================================================
\section{Introduction}
\label{sec:litreview-intro}
% =====================================================================

This chapter situates the present work within the existing literature
on residential, grid-tied rooftop photovoltaic (PV) systems, with a
geographic emphasis on Egypt. The aim is twofold: first, to summarise
the methodological state of the art across the four technical pillars
that the thesis combines --- PV energy modelling, rooftop area
estimation, electricity-tariff-aware financial modelling, and
uncertainty quantification --- and second, to surface the specific
gaps that the methodology of Chapter~3 addresses.

Three deliberate scope restrictions apply. The chapter considers
\emph{residential} systems (typically $1$--$15$\,kWp) rather than
utility-scale PV plants, whose financing, siting, and grid-coupling
problems differ qualitatively. It considers \emph{grid-tied} systems
(self-consumption with surplus export) rather than off-grid or hybrid
diesel/battery systems, which dominate the rural-electrification
literature. And it considers conventional roof-mounted modules rather
than building-integrated PV (BIPV) or bifacial trackers. With those
boundaries set, \S\ref{sec:lit-pv}--\S\ref{sec:lit-egypt} review the
literature thematically; \S\ref{sec:lit-gap} synthesises the four
gaps that motivate Chapter~3.

% =====================================================================
\section{Photovoltaic energy modelling and simulation}
\label{sec:lit-pv}
% =====================================================================

\subsection{From first-principles to library toolchains}
\label{sec:lit-pv-tools}

The conversion of broadband solar irradiance into AC electrical
output at the meter follows a well-understood physical chain:
horizontal irradiance is decomposed into beam and diffuse components
\citep{liu1960}, transposed onto the tilted plane of the array, used
together with the solar position \citep{nrel-spa} to determine
plane-of-array (POA) irradiance, combined with module temperature and
temperature coefficients to yield DC output, and finally passed
through inverter and balance-of-system efficiencies to produce AC
energy. The classic textbook treatment is given by
\citet{duffie}; \citet{parida2011} survey the underlying cell-level
technologies and \citet{green2024} provide the most recent
record-efficiency tables.

Several open and semi-open simulation toolchains implement this
pipeline. The U.S.\ Department of Energy's PVWatts \citep{pvwatts2014}
exposes a deliberately simplified model designed for quick economic
screening. The full System Advisor Model (SAM) \citep{nrel-sam}
exposes a richer detailed-loss decomposition and is the de-facto
reference for utility-scale studies in North America. The
\texttt{pvlib} Python package \citep{pvlib} re-implements the
underlying algorithms as a library suitable for embedding in larger
analytical pipelines, which is the approach this thesis takes. The
European Commission's PVGIS \citep{pvgis} bundles a TMY climate
database with a similar conversion model and exposes both as a public
REST endpoint. \citet{ieapvps-task7} place these tools in the wider
context of PV deployment trends.

\subsection{Cross-validation and dual-model precedent}
\label{sec:lit-pv-cross}

Although the underlying physics is shared, every implementation makes
different choices about loss decomposition, transposition algorithm,
spectral correction, and low-irradiance behaviour
\citep{mavromatakis2010}, and these can produce annual-yield
disagreements of several percent on the same input data.
\citet{freeman2014} validate SAM against measured plant data and
report site-dependent biases; \citet{gilman2018} document the SAM
photovoltaic model in detail and quantify model-induced uncertainty.
\citet{polo2020} benchmark several simulation chains against measured
output across multiple sites and conclude that no single model
dominates --- agreement \emph{between} models is itself useful
evidence of robustness. Despite this, residential-scale PV studies
overwhelmingly report a single-model output as the headline yield
figure. Cross-validating two independent implementations against the
same input data is uncommon outside utility-scale validation reports
and represents the methodological pattern that
\S\ref{sec:pv} of this thesis adopts.

% =====================================================================
\section{Rooftop area estimation from remote data}
\label{sec:lit-roof}
% =====================================================================

A PV yield calculation is only as accurate as the roof area on which
the system is sized. The literature on automatically estimating that
area can be organised by data source.

\subsection{Manual entry, GIS, and cadastral data}
\label{sec:lit-roof-manual}

The default in commercial homeowner-facing calculators (PVWatts,
EnergySage, vendor sizing tools) is to ask the user to enter their
roof area or, equivalently, the number of panels they think will
fit. This is convenient but unreliable: homeowners systematically
mis-estimate roof dimensions, and the resulting yield projections
inherit the error linearly. Cadastral and GIS approaches replace
manual entry with authoritative geometry. \citet{wiginton2010} use
LiDAR returns and municipal building polygons to estimate Ontario's
aggregate rooftop solar potential; \citet{brito2012} apply LiDAR-based
extraction in a Lisbon suburb; \citet{mainzer2014} build a
high-resolution GIS pipeline for German rooftops; and
\citet{lukac2014} rasterise LiDAR returns to derive per-roof
suitability scores. These approaches work well where high-quality
LiDAR and cadastral records are publicly available --- a condition
that is not met in Egypt or most of North Africa, where municipal
geometry is sparse or undigitised.

\subsection{Satellite imagery: classical CV and deep learning}
\label{sec:lit-roof-cv}

Where authoritative geometry is unavailable, satellite or aerial
imagery becomes the data source. Classical computer-vision pipelines
combine edge detection, region growing, and shape regularisation;
\citet{wang2018} extract building footprints from very-high-resolution
imagery using a deep network followed by a guided-filter
post-processing stage, illustrating how classical regularisation
remains useful even alongside learned segmentation. Pure
deep-learning approaches now dominate the segmentation literature.
\citet{ronneberger2015} introduced the U-Net architecture that
underpins most modern building-segmentation models;
\citet{assouline2017} apply machine learning to large-scale rooftop
solar potential estimation in Switzerland; \citet{yu2018deepsolar}
build the influential \emph{DeepSolar} database covering the
contiguous United States; and \citet{lee2019deeproof} extend this to
solar potential rather than mere footprint extraction.

The recurring caveat is training data. \citet{vargas-munoz2021}
document the substantial label noise and geographic bias that
characterise OpenStreetMap (OSM) and other crowdsourced annotations
when they are used as training targets. North African rooftops are
under-represented in the corpora behind DeepSolar and DeepRoof, and
no comparable annotated dataset exists for Egyptian residential
rooftops. A purely deep-learning approach therefore inherits a
training-distribution mismatch that is not straightforward to remedy
within the scope of a residential feasibility tool.

\subsection{OpenStreetMap-based footprints}
\label{sec:lit-roof-osm}

A pragmatic alternative is to consume OSM building footprints
directly via the Overpass API \citep{osm-overpass}. \citet{osm-quality}
benchmark OSM geometry against authoritative survey data and find
quality varies but is workable in built-up areas; \citet{vargas-munoz2021}
discuss the same tradeoff in the machine-learning context. The OSM
route avoids both the cadastral coverage gap and the training-data
mismatch, but raw OSM polygons are noisy: vertex counts vary,
right-angle constraints are violated, and small jitter accumulates in
area calculations. Combining OSM footprints with classical
computer-vision regularisation against satellite imagery --- using
imagery to refine geometry rather than to learn it from scratch --- is
notably underexplored in the rooftop-PV literature.

% =====================================================================
\section{Electricity tariffs and PV economic modelling}
\label{sec:lit-tariff}
% =====================================================================

\subsection{Tariff structures: flat, time-of-use, and tiered}
\label{sec:lit-tariff-structures}

Residential electricity tariffs fall into three broad families
\citep{bhattacharyya2019}. \emph{Flat} tariffs charge a single
$\mathrm{price/kWh}$ regardless of consumption or time;
\emph{time-of-use} (ToU) tariffs vary the price across diurnal blocks;
and \emph{increasing-block} or \emph{tiered} tariffs apply a
piecewise-increasing schedule of marginal prices to consumption
within consumption bands. \citet{borenstein2012} analyses the
distributional consequences of nonlinear (tiered) pricing and shows
that a kWh saved in the highest block is worth several times as much
as a kWh saved in the lowest block. \citet{darghouth2011} demonstrate
in the Californian context that the bill savings produced by
distributed PV depend strongly on the tariff schedule: under a
flat tariff PV is essentially indifferent to load shape, but under a
tiered or ToU tariff the marginal value of the avoided kWh depends on
which block it would otherwise have fallen in.

The implication for PV calculators is direct. Tools such as PVWatts
\citep{pvwatts2014} and the simplified models surveyed by
\citet{audenaert2010} compute bill savings against a single average
tariff, which is correct only for flat-rate consumers. Where the
underlying tariff is tiered, this approach systematically
under-estimates savings for high-consumption households (whose
displaced kWh come from expensive top blocks) and over-estimates them
for low-consumption households (whose displaced kWh come from cheaper
blocks).

\subsection{Financial metrics: payback, NPV, and LCOE}
\label{sec:lit-tariff-metrics}

Once an annual bill saving is computed, several headline financial
metrics summarise the investment. Simple payback period is the
ratio of capital expenditure to first-year saving and is favoured for
its homeowner-legibility. Discounted payback corrects for the time
value of money. Net present value (NPV) sums the discounted
twenty-five-year cash-flow stream. Levelised cost of electricity
(LCOE) divides total lifetime cost by total lifetime generation;
\citet{branker2011} review LCOE methodology specifically for solar PV
and \citet{talavera2010} provide a comprehensive sensitivity analysis
of the internal rate of return for grid-connected systems.
\citet{irena2023} report cost-trend evidence that residential LCOE has
fallen below grid retail tariffs across most of the world.

The economics also depend on the export rule: \emph{net metering}
credits exported kWh at the retail tariff (treating the meter as
running backwards), \emph{net billing} or self-consumption-first
schemes credit exports at a lower wholesale or feed-in rate, and
\emph{feed-in tariffs (FIT)} pay an explicit per-kWh price for all
generation. These distinctions interact non-trivially with tiered
retail tariffs, since the value of a self-consumed kWh is the
\emph{marginal} retail price avoided, which under a tiered schedule
is itself a function of household consumption.

\subsection{Egypt-specific context: the EgyptERA progressive tariff}
\label{sec:lit-tariff-egypt}

Egypt's residential electricity sector operates under Electricity Law
87/2015 \citep{egypt-electricity-act}, which established the
regulator (EgyptERA) and provided the framework for the present
seven-tier increasing-block residential schedule
\citep{egyptera_tariff}. Marginal prices in the highest tier are
several multiples of those in the lowest tier, and the schedule has
been revised upwards repeatedly since the post-2014 reform path. Grid
emissions are reported by EEHC \citep{eehc_emissions} and have
declined as natural gas has displaced fuel oil in generation.

The Egyptian feasibility literature has not, to date, modelled this
tariff structure faithfully. \citet{mahmoud2023} present a
pre-feasibility analysis of residential rooftop PV in Egypt under
average-rate assumptions; \citet{esmail-negm-2021} carry out a
techno-economic analysis of grid-tied PV using a single representative
tariff figure; \citet{mohamed2020} treat a stand-alone hybrid
configuration (off-scope for this thesis but illustrative of the
methodological gap). In each case a flat-rate assumption is either
made explicitly or implied by the use of an average $\mathrm{EGP/kWh}$
input. The marginal-block sensitivity that
\citet{borenstein2012} and \citet{darghouth2011} identify as the
dominant determinant of household-PV economics is therefore absent
from the published Egyptian analyses.

% =====================================================================
\section{Uncertainty quantification in PV finance}
\label{sec:lit-uncertainty}
% =====================================================================

\subsection{Why point estimates mislead homeowners}
\label{sec:lit-uncertainty-point}

A residential PV business case depends on at least seven uncertain
parameters: module degradation rate \citep{jordan2013}, inverter
service life \citep{ieapvps2021}, capital expenditure (subject to
import-cost and exchange-rate risk in the Egyptian context),
ongoing operations and maintenance, soiling losses
\citep{elsayed2017}, retail tariff escalation, and weather-driven
year-to-year yield variability. A point-estimate payback figure
collapses all of these into a single number. \citet{spertino2014}
analyse long-term rooftop PV economics across the two main European
markets and show that headline payback is highly sensitive to the
chosen central values for these parameters; small changes in the
discount rate or tariff-escalation assumption can move the payback by
years.

\subsection{Sensitivity analysis and Monte Carlo precedent}
\label{sec:lit-uncertainty-mc}

Two complementary techniques address this problem. One-at-a-time
(OAT) sensitivity analysis, summarised by \citet{saltelli2008},
perturbs each input across a plausible range with all others held at
their central value, producing a tornado chart that ranks parameters
by their leverage on the output. Monte Carlo (MC) simulation samples
all uncertain parameters jointly from probability distributions and
reports the resulting output distribution as percentile bands.
\citet{drury2012} apply MC sampling within the U.S.\ SunShot Vision
Study to bound rooftop-PV market projections; \citet{lacirignola2017}
combine sensitivity and MC analyses for emerging PV technologies
within a life-cycle assessment framework; \citet{cucchiella2012}
apply MC to building-integrated PV economics. The IEA-PVPS Task 12
report \citep{ieapvps2020} provides standardised guidelines for
LCA-style uncertainty propagation; the IPCC \citep{ipcc-ar6-wg3}
similarly motivates probabilistic reporting of mitigation outcomes,
and the EPA \citep{epa-equivalencies} and EEA
\citep{eea-passenger-cars} provide the equivalence factors used to
translate avoided $\COtwo$ into homeowner-legible units. Yet most of
this MC literature targets either utility-scale finance or LCA-style
environmental indicators rather than residential payback
distributions reported back to a homeowner. The typical homeowner
deliverable in published Egyptian and MENA studies remains a single
number.

% =====================================================================
\section{Egypt-specific residential PV studies}
\label{sec:lit-egypt}
% =====================================================================

The technical solar resource in Egypt is exceptional: the
NREA Solar Atlas \citep{nrea-atlas} and Solargis assessments
\citep{solargis-egypt} report annual GHI of $2{,}000$--$2{,}600$\,
kWh\,m$^{-2}$ across most of the country, with Cairo and Upper Egypt
in the upper end of that range. \citet{khalil2018} compare statistical
and inclined-surface irradiance models against Egyptian measurements;
\citet{elsayed2017} characterise soiling losses under Cairo's
atmospheric conditions and report that uncleaned modules can lose
several percent per month during dust season --- a loss term that
materially shifts payback in the desert climate.

The residential techno-economic literature for Egypt remains sparse.
\citet{mahmoud2023} and \citet{esmail-negm-2021} are the two most
directly relevant studies. Both establish that grid-tied rooftop PV
is economically attractive in Egypt against current tariffs; both
adopt a deterministic single-system framing; and both treat the
tariff as flat. Neither incorporates automated roof geometry, tiered
billing, or probabilistic uncertainty bounds. Wider regional context
is provided by \citet{mohamed2020} for stand-alone systems and by
the IEA's renewables forecasts \citep{iea-renewables-2023} for
aggregate market trajectories.

Two technology areas adjacent to this thesis but not central to it
deserve mention for completeness. Maximum-power-point tracking and
inverter design have a mature literature: \citet{esram2007} survey
MPPT algorithms and \citet{kjaer2005} review single-phase
grid-connected inverter topologies. The present work treats inverter
behaviour as a fixed efficiency curve and inherits MPPT performance
from the manufacturer specification, so these references serve only
to bound the abstraction.

% =====================================================================
\section{Synthesis: research gap and contributions}
\label{sec:lit-gap}
% =====================================================================

The four preceding sections each end on the same note: a methodological
choice that is well-supported in the wider PV literature but absent
from the published Egyptian residential analyses. Taken together they
define the gap that this thesis addresses. Table~\ref{tab:gap-map}
summarises the mapping from each gap to the corresponding
methodological contribution in Chapter~3.

\begin{table}[h]
  \centering
  \small
  \begin{tabularx}{\textwidth}{@{} l X X @{}}
    \toprule
    \textbf{\#} & \textbf{Gap identified in the literature}
                & \textbf{This thesis (Chapter~3)} \\
    \midrule
    1 & Residential PV studies report a single-model annual yield;
        cross-validation between independent implementations is rare
        outside utility-scale validation reports
        \citep{freeman2014, polo2020}.
      & Dual independent energy models (\texttt{pvlib} library
        chain~+ first-principles physics chain) producing two
        yield estimates whose agreement is itself reported as
        evidence of robustness
        --- \S\ref{sec:pv}. \\
    \addlinespace
    2 & Roof area is either entered manually (unreliable) or
        extracted by deep-learning models trained on corpora that
        under-represent North African rooftops
        \citep{vargas-munoz2021, yu2018deepsolar}; LiDAR-cadastral
        approaches \citep{wiginton2010, brito2012} require data
        unavailable in Egypt.
      & OSM building footprints retrieved through Overpass and
        regularised with classical computer-vision techniques against
        satellite imagery, avoiding both manual entry and
        training-distribution mismatch
        --- \S\ref{sec:roof}. \\
    \addlinespace
    3 & Egyptian residential PV studies
        \citep{mahmoud2023, esmail-negm-2021} treat the EgyptERA
        seven-tier progressive tariff as a single average rate, even
        though the marginal-block effect is the dominant driver of
        household-PV economics under nonlinear pricing
        \citep{borenstein2012, darghouth2011}.
      & Self-consumption-first netting against the actual EgyptERA
        block schedule \citep{egyptera_tariff}, with system size
        optimised over discrete options to maximise NPV
        --- \S\ref{sec:tariff}. \\
    \addlinespace
    4 & Residential PV economics are reported as point-estimate
        payback figures, even though seven or more parameters are
        independently uncertain \citep{jordan2013, ieapvps2021,
        spertino2014}; published Monte Carlo work targets utility
        finance or LCA, not homeowner-facing decision support
        \citep{drury2012, lacirignola2017}.
      & Monte Carlo simulation over seven stochastic inputs producing
        percentile fan-charts on payback and NPV, alongside an OAT
        sensitivity tornado for parameter attribution
        --- \S\ref{sec:montecarlo}. \\
    \bottomrule
  \end{tabularx}
  \caption{Mapping from gaps identified in the literature to the
  contributions of Chapter~3.}
  \label{tab:gap-map}
\end{table}

The methodology presented in Chapter~3 is structured around these
four contributions and addresses each of the gaps in turn. The
validation chapter that follows it cross-checks the resulting system
against the same Egyptian studies surveyed above
\citep{mahmoud2023, esmail-negm-2021} as well as against published
LCOE and payback ranges \citep{branker2011, irena2023}, closing the
loop between the literature and the implemented system.
